\documentclass[conference]{IEEEtran}

\usepackage{amsmath,amssymb,amsfonts}
\usepackage{graphicx}
\usepackage{textcomp}
\usepackage{xcolor}
\usepackage{hyperref}
\usepackage{booktabs}
\usepackage{multirow}
\usepackage{subcaption}
\usepackage[backend=biber,style=ieee]{biblatex}
\addbibresource{references.bib}

\begin{document}

\title{AudioRestore: Text-Controlled Audio Quality Restoration\\via Fine-Tuned SonicMaster}

\author{
\IEEEauthorblockN{[Your Name]}
\IEEEauthorblockA{CS 614 -- Applications of Machine Learning\\
George Mason University\\
Spring 2026}
}

\maketitle

% ============================================================================
\begin{abstract}
Audio recordings are frequently degraded by lossy compression codecs such as
MP3 and OGG Vorbis, resulting in audible artifacts and loss of high-frequency
content. We present \textbf{AudioRestore}, a proof-of-concept system for
text-controlled audio quality restoration built by fine-tuning SonicMaster, a
0.9-billion-parameter generative model that uses flow matching conditioned on
natural language prompts. We construct a codec-degraded dataset by applying MP3
and OGG compression at various bitrates (24--128\,kbps) to clean music from
MUSDB18-HQ, paired with descriptive text prompts. Fine-tuning the pretrained
SonicMaster model on this dataset specializes it for codec artifact removal,
enabling users to control the restoration process through prompts such as
``remove MP3 artifacts at 64\,kbps.'' We evaluate our system using SDR and
SI-SNR metrics across multiple bitrates and demonstrate improved restoration
quality over the general-purpose pretrained model. Our work shows how a
state-of-the-art generative audio model can be adapted for a specific
real-world application with minimal data and compute.
\end{abstract}

\begin{IEEEkeywords}
audio restoration, codec artifacts, flow matching, SonicMaster, fine-tuning,
text-controlled audio enhancement
\end{IEEEkeywords}

% ============================================================================
\section{Introduction}

Billions of audio files exist in compressed formats. Streaming services like
Spotify Free deliver music at 128\,kbps, YouTube rips circulate at even lower
bitrates, and decades of MP3 collections were encoded with aggressive
compression. Lossy codecs introduce audible artifacts---pre-echo, spectral
holes, and loss of high-frequency detail---that degrade the listening
experience~\cite{li2025apollo}.

Traditional audio restoration relies on manual signal processing chains:
equalization, de-noising, and spectral repair, each requiring expert tuning.
Recent deep learning models such as Apollo~\cite{li2025apollo} and
SonicMaster~\cite{melechovsky2025sonicmaster} automate this process, but
general-purpose models may not optimally handle specific degradation types.

We propose \textbf{AudioRestore}, a system that fine-tunes SonicMaster on
codec-specific degradations. Our contributions are:

\begin{enumerate}
    \item A \textbf{codec-degraded dataset} with paired clean/degraded 30-second
    clips at various MP3 and OGG bitrates, annotated with natural language
    prompts describing the degradation.
    \item A \textbf{fine-tuning procedure} that adapts SonicMaster's 0.9B-parameter
    flow-matching model to specialize in codec artifact removal while
    preserving its text-controlled interface.
    \item A \textbf{proof-of-concept demo} showing text-controlled restoration
    with before/after spectrograms and quantitative metrics.
\end{enumerate}

% ============================================================================
\section{Related Work}

\subsection{Audio Super-Resolution and Bandwidth Extension}
Audio super-resolution aims to reconstruct high-frequency content from
bandwidth-limited signals. Methods range from traditional interpolation to
deep learning approaches using CNNs and GANs~\cite{su2021bandwidth}. These
methods typically operate in a fixed input-output configuration without
user control.

\subsection{Apollo}
Apollo~\cite{li2025apollo} is a band-sequence model for music restoration in
compressed audio, published at ICASSP 2025. It splits the spectrogram into 80
frequency bands, processes them through Roformer attention layers with
inter-band communication blocks (ICB), and reconstructs the waveform via a GAN
framework. Apollo achieves strong results on MUSDB18-HQ but operates without
any text conditioning---it is a fixed input-to-output mapping.

\subsection{SonicMaster}
SonicMaster~\cite{melechovsky2025sonicmaster} is a unified generative model
for music restoration and mastering, submitted to ICLR 2026. It addresses five
categories of audio artifacts (equalization, dynamics, reverb, amplitude,
stereo) through a single flow-matching model conditioned on natural language
prompts. The architecture uses:
\begin{itemize}
    \item A \textbf{Stable Audio VAE}~\cite{evans2024stable} encoder/decoder to
    convert between waveforms and latent representations.
    \item A \textbf{TangoFlux transformer}~\cite{ghosh2024tangoflux} ($\sim$515M
    parameters) that performs flow-matching in latent space, conditioned on
    text embeddings from a T5 encoder and audio condition embeddings.
    \item \textbf{Text and audio conditioning} through pooled projections that
    guide the restoration direction.
\end{itemize}
SonicMaster is trained on the SonicMasterDataset (25k paired tracks, 208
hours) with 19 degradation functions. Our work fine-tunes this pretrained model
specifically for codec artifact removal.

% ============================================================================
\section{Method}

\subsection{Architecture}
We use the full SonicMaster architecture without modification. The model
processes 30-second stereo audio chunks at 44.1\,kHz. The degraded input is
encoded into a latent representation by the Stable Audio VAE, transformed by
the TangoFlux flow-matching model conditioned on a text prompt, and decoded
back to a waveform.

% TODO: Include architecture diagram from the plan
% \begin{figure}[t]
%     \centering
%     \includegraphics[width=\linewidth]{figures/architecture.pdf}
%     \caption{SonicMaster architecture for audio restoration.}
%     \label{fig:arch}
% \end{figure}

\subsection{Codec Degradation Pipeline}
Our key modification is a custom degradation pipeline targeting codec
artifacts. We apply:

\begin{itemize}
    \item \textbf{MP3 compression} at bitrates $\in \{24, 32, 48, 64, 96,
    128\}$\,kbps using \texttt{torchaudio.functional.apply\_codec}.
    \item \textbf{OGG Vorbis compression} at quality levels $\in \{-1, 0, 1,
    2, 3, 5\}$.
\end{itemize}

For each clip, we randomly select a codec type (70\% MP3, 30\% OGG) and
parameter, then generate a paired (clean, degraded) sample. Clips are
normalized to 0.95 peak amplitude before degradation.

\subsection{Text Prompt Design}
Each training pair is annotated with two prompts:
\begin{itemize}
    \item A \textbf{specific prompt} describing the degradation, e.g.,
    ``restore audio compressed at 64\,kbps MP3.''
    \item An \textbf{alternative prompt} with generic phrasing, e.g.,
    ``improve the sound quality.''
\end{itemize}

During training, we follow SonicMaster's text augmentation strategy: 10\%
chance of empty prompt (unconditional), 10\% chance of a generic prompt, and
40/40 split between the specific and alternative prompts. This teaches the
model to respond to various phrasings while maintaining unconditional
capability.

\subsection{Fine-Tuning Procedure}
We fine-tune from the pretrained SonicMaster checkpoint (0.9B parameters):
\begin{itemize}
    \item \textbf{Frozen components}: T5 text encoder (parameters fixed
    throughout).
    \item \textbf{Trained components}: TangoFlux transformer, text projection
    layers (\texttt{fc\_text}, \texttt{fc\_text\_audio}), and audio
    conditioning module (\texttt{audio\_cond}).
    \item \textbf{Learning rate}: $1 \times 10^{-5}$ with cosine schedule and
    200-step warmup.
    \item \textbf{Batch size}: 2 per device with 4 gradient accumulation steps
    (effective batch size 8).
    \item \textbf{Precision}: Mixed FP16 for memory efficiency on a single GPU.
    \item \textbf{Epochs}: 30 (early stopping on validation loss).
\end{itemize}

All audio is pre-encoded into VAE latents before training to avoid redundant
VAE forward passes during fine-tuning.

% ============================================================================
\section{Experiments}

\subsection{Dataset}
We construct our dataset from MUSDB18-HQ~\cite{rafii2017musdb18}, extracting
30-second clips from the training split. Each clip is degraded with a randomly
selected codec and bitrate/quality level. We generate approximately 1,000--2,000
paired clips (90\% train, 10\% validation).

\subsection{Evaluation Metrics}
We evaluate using two standard audio quality metrics~\cite{le2019sdr}:
\begin{itemize}
    \item \textbf{SDR} (Signal-to-Distortion Ratio): Measures overall signal
    fidelity in dB.
    \item \textbf{SI-SNR} (Scale-Invariant Signal-to-Noise Ratio): Measures
    signal quality invariant to scaling.
\end{itemize}

\subsection{Results}

% TODO: Fill in with actual results after running fine-tuning
\begin{table}[t]
\centering
\caption{Restoration quality (SDR in dB) at various MP3 bitrates.}
\label{tab:results}
\begin{tabular}{lcccc}
\toprule
\textbf{Bitrate} & \textbf{Degraded} & \textbf{Base SM} & \textbf{Fine-tuned} & \textbf{Gain} \\
\midrule
32\,kbps  & --  & --  & --  & --  \\
64\,kbps  & --  & --  & --  & --  \\
128\,kbps & --  & --  & --  & --  \\
\bottomrule
\end{tabular}
\vspace{0.5em}
\small SM = SonicMaster. Gain = Fine-tuned SDR $-$ Degraded SDR.
\end{table}

% TODO: Add spectrogram figures
% \begin{figure}[t]
%     \centering
%     \includegraphics[width=\linewidth]{figures/spectrograms_64k.png}
%     \caption{Spectrograms at 64\,kbps MP3. Left: original. Center: degraded.
%     Right: restored by fine-tuned SonicMaster.}
%     \label{fig:spectrograms}
% \end{figure}

% ============================================================================
\section{Discussion}

\subsection{Text Control Advantage}
Unlike Apollo, which provides no user control over the restoration process,
SonicMaster's text conditioning allows users to specify the type of
degradation they want to address. Our fine-tuning preserves this capability
while specializing the model for codec artifacts.

\subsection{Limitations}
\begin{itemize}
    \item Fine-tuning was performed on a relatively small dataset due to
    compute constraints.
    \item The model processes 30-second chunks; longer files require chunked
    processing with crossfade stitching.
    \item Inference requires a GPU with at least 8\,GB VRAM (FP16).
\end{itemize}

% ============================================================================
\section{Product Vision}

\textbf{AudioRestore} could be deployed as a consumer service:
\begin{itemize}
    \item A web app or API where users upload compressed audio and describe the
    desired enhancement in natural language.
    \item Integration with music streaming platforms to upscale low-bitrate
    streams in real-time.
    \item A plugin for digital audio workstations (DAWs) for mastering
    engineers working with degraded source material.
    \item A mobile app for restoring old MP3 collections or voice recordings.
\end{itemize}

The text-controlled interface makes it accessible to non-experts who cannot
operate traditional audio restoration tools.

% ============================================================================
\section{Conclusion}

We presented AudioRestore, a proof-of-concept system for text-controlled audio
quality restoration. By fine-tuning SonicMaster on codec-degraded audio with
custom text prompts, we demonstrated that a general-purpose audio restoration
model can be specialized for a specific degradation type with minimal data and
compute. The system allows users to control the restoration process through
natural language, making high-quality audio restoration accessible to
non-experts. Future work includes expanding the degradation pipeline, training
on larger datasets, and optimizing for real-time inference.

% ============================================================================
\printbibliography

\end{document}
